% !TEX root = ../paper.tex
\doclearpage
\section{Introduction}

With the rise of large language models~\cite{OuyangWJAWMZASRSHKMSAWCLL22,TouvronLIMLLRGHARJGL23}, text has become a dominant form of data, and analytics workloads reflect it. Across more than $60{,}000$ real-world business-intelligence repositories, $48\%$ of stored values are strings~\cite{VogelsgesangHFKLMNT18}, and across the Amazon Redshift fleet, strings make up $52.1\%$ of all columns~\cite{RenenHPVDNLSK24}. Even $6$ of the $22$ TPC-H queries~\cite{PoessF00} filter rows with \texttt{LIKE} patterns.

Much of this text-heavy data lives in cloud warehouses such as Redshift, on infrastructure the data owner does not control, so users must simply trust the server to evaluate their queries, string predicates included, correctly. Verifiable databases remove this trust: the server returns, alongside every query result, a succinct proof that the result is correct~\cite{ZhangKC15,ZhangGKPP17,LiWXWR23,GuFN25,SXT,qedb,cryptoeprint:2026/1260}.

\parhead{Research Gap}Surprisingly, none of these systems supports operations on the contents of strings. Across design paradigms, from authenticated data structures~\cite{ZhangKC15,qedb} to circuit-based~\cite{LiWXWR23,GuFN25} and interactive-proof-based systems~\cite{ZhangGKPP17,SXT,cryptoeprint:2026/1260}, string columns are treated as opaque values: each string is encoded or hashed into a single field element, and equality is the only supported string operation. Proof of SQL, for instance, states that it does not support any string operation beyond $=$ and $\neq$~\cite{posql}. The omission shows in every evaluation. vSQL discards the columns used by substring and wildcard predicates and leaves such queries to future work~\cite{ZhangGKPP17}; ZKSQL omits the pattern-matching predicate of TPC-H Q9 as beyond its scope~\cite{LiWXWR23}; PoneglyphDB follows ZKSQL and excludes it as well~\cite{GuFN25}; and $\truthtable$ removes every string operation other than equality from the TPC-H queries it verifies~\cite{cryptoeprint:2026/1260}. Strikingly, this includes the circuit-based systems, which could in principle embed the characters of every string directly in the circuit.

We believe this gap reflects a genuine difficulty rather than an oversight. Proving a white-box string operation requires reasoning about the individual characters of every string and their positions, for example that the factors of a pattern occur in order and without overlap, rather than about a single value per row. Hashing a string into a field element destroys exactly this structure, while exposing it naively multiplies the size of the proven statement by the length of the strings. Prior work does prove pattern matching over a single committed document~\cite{PapadopoulosPTT15,AngelIMSW24}, but not over the columns of a database, where strings have varying lengths, are filtered, joined and rematerialized by the query plan, and must remain consistent with the rest of the table throughout. As a result, a large and growing class of real-world queries lies outside the reach of verifiable databases, and this gap has so far been left open.

\ali{Continue: our approach (string arithmetization with auxiliary columns, string toolbox, wild-card pattern matching, fingerprint pre-filtering) and the list of contributions.}
